% MOCO: Measure Once, Compile Once. BreaQ 2026, Quantum AI track, Task 2.
% Every number in this file is copied from a tracked CSV under results/ or from results/task2/TABLES.md;
% the source is given in a comment next to it. Means are over 5 seeds; "±" is the sample standard deviation.
% The test set is used only for reporting.
%
% Citation keys (interim bibliography paper/refs_verified.bib, every entry verified via Crossref or arXiv):
%   dinut2026 dalpozzolo2015 schuld2019gradients cerezo2021vqa mcclean2018barren grant2019init wang2021noise
%   perezsalinas2020reuploading schuld2021encoding schreiber2023surrogates huang2020shadows
%   cerezo2023simulability bowles2024benchmark havlicek2019 hoeffding1963 maurer2009 waudbysmith2024betting
\documentclass[10pt,twocolumn]{article}
\usepackage[margin=0.8in,columnsep=0.25in]{geometry}
\usepackage[T1]{fontenc}
\usepackage{lmodern}
\usepackage{amsmath,amssymb,amsthm}
\usepackage{graphicx}
\usepackage{booktabs}
\usepackage{microtype}
\usepackage[numbers,sort&compress]{natbib}
\usepackage{xcolor}
\usepackage[font=small,labelfont=bf]{caption}
\usepackage{enumitem}
\usepackage[colorlinks=true,linkcolor=blue!50!black,citecolor=blue!50!black,urlcolor=blue!50!black]{hyperref}
\usepackage[capitalise,noabbrev]{cleveref}
\setlist{nosep,leftmargin=*}

\newtheorem{theorem}{Theorem}
\newtheorem{corollary}{Corollary}
\theoremstyle{definition}
\newtheorem{definition}{Definition}
\newcommand{\figdir}{../results/task2/figures}
\newcommand{\ex}[1]{\langle #1\rangle}

\title{MOCO: Measure Once, Compile Once\\[2pt]
\large A noise-invariant variational classifier for credit-card fraud, trained with two circuit executions per example}
\author{Darius Tatomir \and Deliu Georgiana\\[2pt]\small BreaQ 2026, Quantum AI track, Task 2 ``Fraud on a Quantum Budget''}
\date{}

\begin{document}
\maketitle

\begin{abstract}
Training a variational quantum classifier (VQC) on hardware is expensive: every gradient step costs circuit
executions, and depolarizing noise both compresses the scores and reorders them. We present MOCO (Measure Once,
Compile Once), a four-qubit product circuit, RY encoding followed by one trainable rotation per qubit and a weighted
$Z$ readout. In the Heisenberg picture the readout is linear in the Bloch vectors of the encoded data, so training
reduces to a convex logistic regression on data measured once (two executions per training transaction); the
optimum is compiled back into rotation angles and observable weights. Because every qubit sees the same number of
gates and no entangling gate, depolarizing noise multiplies every score's distance to $1/2$ by the same factor
$\lambda^2$, $\lambda = 1-4p/3$: AUPRC, AUROC and the ranking are identical at every noise level, and the decision
threshold moves analytically. Under the organizers' protocol (common split, 5 seeds, noise after every native
gate), MOCO reaches a test AUPRC of $0.716\pm0.007$ at every $p$, while the kit baseline falls from
$0.580\pm0.053$ to $0.405\pm0.177$ at $p=0.05$; a logistic regression on the same inputs reaches
$0.710\pm0.007$, so MOCO matches, and does not beat, the classical reference. Training and model selection cost
6,490 executions instead of 2,859,360. With finite shots the noise becomes a shot tax of $\lambda^{-4}$, which an
anytime-valid adaptive triage reduces to $208.5$ shots per transaction at $p=0.05$ (1024 fixed) at the same F1.
The circuit is classically simulable and we claim no quantum advantage.
% Numbers: moco_per_seed_single_10.csv, baseline_per_seed.csv, ledger_final.csv, ledger_moco_single_10.csv,
% ledger_baseline.csv, adaptive_shots.csv (MOCO kit, adaptive, cap 4096, p = 0.05, mean shots 208.489).
\end{abstract}

\section{Introduction}

Variational quantum algorithms \citep{cerezo2021vqa} train parametrized circuits with a classical optimizer. On
hardware, each gradient is estimated from circuit executions: with the parameter-shift rule
\citep{schuld2019gradients}, a circuit with $P$ trainable angles costs $1+2P$ executions per example and per
epoch. The kit's baseline VQC for this task, which follows the pipeline of \citet{dinut2026}, is a ZZ feature map
followed by an EfficientSU2-style ansatz with 32 parameters trained for 30 epochs on a balanced set of 590 rows;
by the kit's formula it needs 1,150,500 training executions and 1,708,860 validation executions.
% ledger_baseline.csv: train executions 1150500 (= 30 x 590 x 65), validation executions 1708860 (= 30 x 56962).
Noise adds a second problem. Depolarizing channels after each of its 126 noisy gate locations compress the
scores towards $1/2$ non-uniformly, so the ranking of transactions changes as well
\citep{dinut2026,wang2021noise}.
% circuits.csv: baseline VQC (kit) noise_locations 126.

Fraud detection makes the ranking the quantity of interest: the ULB credit-card dataset has 492 frauds among
284,807 transactions \citep{dalpozzolo2015}, and the protocol ranks methods by the area under the
precision-recall curve (AUPRC). A noise model that only compresses the scores uniformly would leave AUPRC
unchanged; one that reorders them cannot be repaired by any number of shots.
% dataset_split.csv: full dataset rows 284807, frauds 492.

\paragraph{Contributions.}
\begin{enumerate}
\item \textbf{MOCO}, a product circuit whose training is an exact convex reparametrization of the VQC loss: the
training data are measured once (two executions per row, 6,490 in total) and the global optimum is compiled back
into the circuit, with no circuit inside the optimization loop.
% ledger_moco_single_10.csv: train executions 6490, executions_per_example 2.
\item A \textbf{rank-invariance theorem} under the official noise model, with an analytic threshold transfer
and a shot tax of exactly $\lambda^{-4}$ in signal compression, checked against the kit's density-matrix
simulator.
\item An evaluation under the \textbf{official protocol} (5 seeds), a cross-validation on all 492 frauds, a
cost ledger under three definitions, ablations and a shallow entangled warm-start variant.
\item \textbf{Adaptive shot triage} with anytime-valid confidence intervals and a \textbf{noise meter} that
estimates $p$ from 1,000 calibration executions, so that the threshold needs no noisy validation pass.
% noise_meter.csv: calibration executions 1000 (MOCO kit, least squares, N_cal = 1000, S = 1024).
\end{enumerate}
We state the limits plainly: MOCO is classically simulable and, as a model class, a logistic regression on the
same angles; its contribution is robustness, stability and cost, not accuracy beyond the classical reference.

\section{Setup}

\paragraph{Data and split.} The dataset has 30 columns: 28 anonymized principal components V1--V28, Time and
Amount \citep{dalpozzolo2015}. The common split is stratified 60/20/20 with seed 42, before any balancing: 170,883
training rows (295 frauds), 56,962 validation rows (99 frauds) and 56,962 test rows (98 frauds). The test fraud
rate, i.e.\ the AUPRC of a random ranking, is 0.172\%.
% dataset_split.csv: train 170883 / 295, validation 56962 / 99, test 56962 / 98, test fraud_rate 0.001720.

\paragraph{Official protocol.} The organizers' evaluation, which we reproduce cell by cell in
\texttt{qai/protocol.py}, decomposes every circuit into $\{RX, RY, RZ, \mathrm{Rot}, H, \mathrm{CNOT},
\mathrm{CZ}\}$ and inserts a depolarizing channel of probability $p$ after every gate, on each qubit it touches,
simulated exactly with density matrices, for $p\in\{0, 0.001, 0.01, 0.05\}$. The threshold maximizes F1 on the
validation scores computed at the same $p$. We report test AUPRC, AUROC, the Spearman correlation $\rho$ between
noisy and noiseless test scores, F1 and MCC, as mean $\pm$ standard deviation over seeds 0--4. Balancing, encoding,
architecture, epochs, hyperparameters and thresholds are chosen on training or validation data only; \textbf{the
test set is used only for reporting}.

\paragraph{Encoding.} As in the kit, Time and Amount are standardized, the 30 columns are projected onto four
principal components, and each component is scaled to an angle in $[0,\pi]$; all three maps are fitted on the
undersampled training subset only.

\paragraph{Kit baseline and its ledger.} The baseline is the VQC of Task 1: 93 native gates, 33 CNOTs, depth 45
and 126 noise locations. Its kit-formula ledger is 2,859,360 executions for training and validation; counted with
\texttt{qml.Tracker}, a training step costs 41 executions per example instead of 65, because PennyLane does not
shift parameters that cannot reach the measured qubit, which gives 2,434,560.
% circuits.csv: baseline VQC (kit) 93 / 33 / 126 / 45. ledger_baseline.csv: total executions 2859360,
% executions_per_example 65, tracker_per_example 41, total executions (Tracker) 2434560.
Removing the 15 gates outside the light cone of the readout leaves the scores unchanged (unit test in
\texttt{tests/test\_protocol.py}), so pruning saves neither noise nor executions.
% circuits.csv: R1 light-cone baseline gates 78 (93 - 78 = 15).

\section{Method}

\paragraph{Circuit.} MOCO uses four qubits. Qubit $q$ receives $RY(x_q)$, which loads angle $x_q$, followed by a
trainable $\mathrm{Rot}(\phi_q,\theta_q,0)$; the readout is $H=\sum_q w_q Z_q$ with $w_q\ge0$, $\sum_q w_q=1$, and
the fraud score is $s(x)=(1-\ex{H})/2$. The circuit has 8 native gates, no CNOT, depth 2 and 8 noise locations.
% circuits.csv: MOCO single 8 / 0 / 8 / 2.

\paragraph{Heisenberg picture.} $RY(x)\lvert0\rangle$ has Bloch vector $r(x)=(\sin x,0,\cos x)$. A single-qubit
rotation $U$ acts on Bloch vectors as a rotation $R_U\in SO(3)$, so
$\ex{Z_q}=\hat z\cdot R_U r(x_q)=m_q\cdot r(x_q)$ with $m_q=R_U^{\top}\hat z$ a unit vector: the trainable rotation only
turns the measurement axis. Hence
\begin{equation}
\ex{H}=\sum_q w_q\, m_q\cdot r(x_q)=\sum_q \big(c_q^{x}\sin x_q + c_q^{z}\cos x_q\big),
\label{eq:linear}
\end{equation}
with $c_q=w_q(m_q^x,m_q^z)$, is linear in the measured features $(\sin x_q,\cos x_q)$.

\paragraph{Measure once, fit convexly.} On hardware, the Bloch components of all four qubits are read in two
settings (X and Z bases), so the training set costs two executions per row, 6,490 for the 3,245 rows of the 1:10
undersampled set. Training the circuit with a logistic loss on~\eqref{eq:linear} plus an intercept is an
$\ell_2$-regularized logistic regression on these 8 measured features: strictly convex, solved to its global
optimum with L-BFGS in PyTorch, with the regularization chosen by 5-fold cross-validation inside the training
set, so that no validation executions are spent and no circuit runs inside the loop.
% ledger_moco_single_10.csv: train executions 6490 = 2 x 3245.

\paragraph{Compile.} Let $\beta_q\in\mathbb R^2$ be the fitted coefficients of qubit $q$ on
$(\sin x_q,\cos x_q)$. Setting $m_q=-\beta_q/\lVert\beta_q\rVert$ (realized by $\mathrm{Rot}(-\phi_q,-\theta_q,0)$,
with $(\theta_q,\phi_q)$ the polar angles of $m_q$) and $w_q=\lVert\beta_q\rVert/\sum_{q'}\lVert\beta_{q'}\rVert$
gives $s(x)=\tfrac12+\sum_q\beta_q\cdot(\sin x_q,\cos x_q)/(2\sum_{q'}\lVert\beta_{q'}\rVert)$, a strictly
increasing affine function of the fitted logit without intercept. The compiled circuit therefore ranks
transactions exactly as the convex model, and the intercept is absorbed by the threshold. The circuit, not a
classical post-processor, performs the classification.

\paragraph{Choices made on validation.} The undersampling ratio and the feature map were chosen on validation
AUPRC over 5 seeds before any test evaluation: single 1:10 reached $0.675\pm0.012$ with 6,490 training
executions; pair 1:30 tied ($0.675\pm0.010$) at 27,435 executions, and the cheaper configuration was kept.
% selection_on_validation.csv (fm, ratio; mean and sample std of val_auprc over seeds 0-4; train_exec).

\section{Theory}

\begin{theorem}[Rank invariance]\label{thm:rank}
Let a circuit on $n$ qubits consist only of single-qubit gates, with exactly $k$ native gates on every qubit,
followed by measurement of $H=\sum_q w_q Z_q$. Under the official noise model, a depolarizing channel
$\mathcal D_p(\rho)=(1-p)\rho+\tfrac p3(X\rho X+Y\rho Y+Z\rho Z)$ after every gate, for every input $x$
\begin{equation}
s_p(x)=\tfrac12-\lambda^{k}\big(\tfrac12-s_0(x)\big),\qquad \lambda=1-\tfrac{4p}{3}.
\label{eq:contraction}
\end{equation}
\end{theorem}

\begin{proof}
A single-qubit state is $\rho=\tfrac12(I+r\cdot\sigma)$ with Bloch vector $r$, $\lVert r\rVert\le1$.
\emph{(i) Channel.} Since $X\sigma_jX=\pm\sigma_j$ with a minus sign exactly for $j\ne x$, and similarly for $Y$
and $Z$, each Pauli component $\sigma_j$ is kept by one of the three Pauli conjugations and flipped by the other
two. Hence $\mathcal D_p$ maps $\sigma_j\mapsto(1-p)\sigma_j+\tfrac p3(1-2)\sigma_j=(1-\tfrac{4p}3)\sigma_j$ and
leaves $I$ fixed: $r\mapsto\lambda r$.
\emph{(ii) Gates.} A unitary $U$ maps $r\mapsto R_U r$ with $R_U\in SO(3)$. Since $\lambda r$ is a scalar
multiple, $R_U(\lambda r)=\lambda R_U r$: the channel commutes with every single-qubit gate.
\emph{(iii) Product structure.} Without multi-qubit gates and with product input $\lvert0\rangle^{\otimes n}$,
every gate and every channel acts on one qubit, so the state remains a product $\bigotimes_q\rho_q$ and each
qubit evolves independently. After its $k$ gates, each followed by a channel, qubit $q$ has Bloch vector
$\lambda^k R_{U_k}\cdots R_{U_1}r_{q,0}=\lambda^k r_q^{\mathrm{ideal}}$ by (i) and (ii).
\emph{(iv) Readout.} $\ex{Z_q}_p=\hat z\cdot\lambda^k r_q^{\mathrm{ideal}}=\lambda^k\ex{Z_q}_0$, so
$\ex{H}_p=\lambda^k\ex{H}_0$ by linearity and $s_p=(1-\lambda^k\ex{H}_0)/2$, which is~\eqref{eq:contraction}.
\end{proof}

For MOCO, $k=2$ ($RY$ and $\mathrm{Rot}$ are both native). Since $\lambda^2>0$ for every $p\ne3/4$, the map
$s_0\mapsto s_p$ is strictly increasing, which gives:

\begin{corollary}[Metrics and threshold]\label{cor:metrics}
For every $p\in[0,3/4)$: (a) AUPRC and AUROC of $s_p$ equal those of $s_0$, and the Spearman correlation between
$s_p$ and $s_0$ is 1; (b) the classifications $\{s_p\ge g_p(t)\}$ and $\{s_0\ge t\}$ coincide with
$g_p(t)=\tfrac12-\lambda^2(\tfrac12-t)$, so the threshold maximizing F1 on noisy validation scores is
$t_p=g_p(t_0)$, and precision, recall, F1 and MCC at the transferred threshold are identical to those at $p=0$.
\end{corollary}

Part (a) needs only that ranking metrics are invariant under strictly increasing maps; part (b) holds because
$g_p$ is a bijection of the score range. The threshold $t_0$ can be computed classically (the noiseless scores of
a product circuit are closed-form), so MOCO needs no noisy validation pass once $p$ is known.

\paragraph{Numerical check.} Against the kit's density-matrix simulator (\texttt{noisy\_predict}), the
contraction~\eqref{eq:contraction} holds to $5.7\times10^{-14}$ on 200 rows at $p=0.05$ and $p=0.5$, the joint
16-outcome distribution equals the product of the per-qubit marginals to $8.6\times10^{-14}$, and the measured
median compression equals $\lambda^2$ to $2.8\times10^{-17}$ at every tested $p$. In the adaptive-shots run, the
transferred threshold matches the threshold chosen on the noisy validation set to $10^{-9}$ or better.
% noise_shots_checks.csv: MOCO "binomial means vs kit noisy_predict" max_error 5.69e-14; "joint probs vs product
% of marginals" 8.55e-14; "compression == lambda^2 (all p)" 2.78e-17. adaptive_checks.csv: "transfer_threshold vs
% max-F1 threshold on noisy validation" rows, limit 1e-9, all PASS.

\paragraph{Shot tax.} On hardware each score is estimated from $S$ shots; the per-shot estimator
$h=\sum_q w_q z_q\in[-1,1]$ has variance at most 1, so the standard error of $\hat s$ is at most
$1/(2\sqrt S)$ at every $p$. By~\eqref{eq:contraction}, the difference between the scores of two transactions
shrinks by $\lambda^2$. Keeping the same signal-to-noise ratio as without noise therefore requires
$S_p/S_0=\lambda^{-4}$ under this bound, i.e.\ $1.32\times$ at $p=0.05$, $3.46\times$ at $p=0.2$ and $81\times$
at $p=0.5$. We define the shot tax of any model as $1/c^2$, with $c$ its measured median compression
$\mathrm{median}\,|s_p-\tfrac12|/|s_0-\tfrac12|$; for MOCO $c=\lambda^2$ exactly. Because the exact per-shot
variance also grows slightly with $p$, the tax is a lower bound on the shots needed.
% shot_tax.csv: MOCO shot_tax 1.3178 (p = 0.05), 3.4578 (p = 0.2), 81.0 (p = 0.5).

\paragraph{Link to Task 1.} The theorem combines two exercises of the Task 1 notebook. Exercise 3.2 shrinks the
scores of a classifier towards 0.5 and observes that AUPRC is unchanged while the recall at a fixed threshold
changes; exercise 10.1 checks that a depolarizing channel after $RY(\theta)$ gives
$\ex Z=(1-\tfrac{4p}3)\cos\theta$. MOCO is the circuit for which the second statement implies the first for the
whole model.

\section{Experiments}

All experiments use seeds 0--4 and the official protocol unless stated otherwise; code and per-seed results are in
the repository (\texttt{results/task2/}).

\subsection{Official protocol}

\begin{table*}[t]
\centering\small
\caption{Test AUPRC under the official protocol (mean $\pm$ std over 5 seeds), Spearman $\rho$ between noisy and
noiseless test scores and F1 at $p=0.05$. LogReg = logistic regression on the same inputs as the model named.}
\label{tab:official}
% Sources: results/task2/TABLES.md (Full protocol tables), from baseline_per_seed.csv, entangled_per_seed.csv,
% moco_per_seed_single_10.csv, moco_per_seed_single_10_q0.5.csv.
\resizebox{\textwidth}{!}{%
\begin{tabular}{lcccccc}
\toprule
Model & $p=0$ & $p=0.001$ & $p=0.01$ & $p=0.05$ & $\rho$ ($p=0.05$) & F1 ($p=0.05$)\\
\midrule
Kit baseline VQC (93 gates, 33 CNOT) & $0.580\pm0.053$ & $0.580\pm0.053$ & $0.564\pm0.049$ & $0.405\pm0.177$ & $0.686\pm0.145$ & $0.537\pm0.184$\\
R2 entangled, random init (15 gates) & $0.635\pm0.018$ & $0.635\pm0.018$ & $0.636\pm0.017$ & $0.638\pm0.015$ & $0.999\pm0.001$ & $0.760\pm0.021$\\
R2 entangled, MOCO warm start & $0.667\pm0.010$ & $0.667\pm0.010$ & $0.667\pm0.010$ & $0.668\pm0.009$ & $0.999\pm0.001$ & $0.798\pm0.009$\\
\textbf{MOCO (8 gates, 0 CNOT)} & $\mathbf{0.716\pm0.007}$ & $\mathbf{0.716\pm0.007}$ & $\mathbf{0.716\pm0.007}$ & $\mathbf{0.716\pm0.007}$ & $\mathbf{1.000}$ & $0.800\pm0.019$\\
LogReg 1:10 (inputs of MOCO) & $0.710\pm0.007$ & = & = & = & -- & $0.794\pm0.023$\\
LogReg 1:1 (inputs of the baseline) & $0.703\pm0.010$ & = & = & = & -- & $0.791\pm0.019$\\
\midrule
MOCO, robust encoding $q=0.5\%$ (not selected) & $0.738\pm0.006$ & $0.738\pm0.006$ & $0.738\pm0.006$ & $0.738\pm0.006$ & $1.000$ & $0.783\pm0.024$\\
LogReg 1:10, same robust inputs & $0.731\pm0.004$ & = & = & = & -- & $0.785\pm0.020$\\
\bottomrule
\end{tabular}}
\end{table*}

\Cref{tab:official} reproduces the task statement for the reference models: the baseline reaches $0.580\pm0.053$
and the logistic regression on its inputs $0.703\pm0.010$ \citep{dinut2026}. At $p=0.05$ the baseline loses 30\%
of its AUPRC, its spread over seeds triples and its ranking decorrelates ($\rho=0.686\pm0.145$); its max-F1
threshold collapses to $0.500\pm0.001$. MOCO keeps $0.716\pm0.007$, $\rho=1$ and F1 $0.800\pm0.019$ at every
noise level, as \cref{cor:metrics} predicts, and matches the logistic regression on the same inputs
($0.710\pm0.007$). \Cref{fig:auprc} shows AUPRC against $p$.
% baseline_per_seed.csv: baseline VQC threshold at p = 0.05: 0.500 +- 0.001. Drop 0.405 / 0.580 = -30%.

\begin{figure}[t]
\centering
\includegraphics[width=\columnwidth]{\figdir/compact/a_auprc_protocol.png}
\caption{Test AUPRC under the official protocol (5 seeds). Source: \texttt{a\_auprc\_protocol.png}, from the
per-seed CSVs of \cref{tab:official}.}
\label{fig:auprc}
\end{figure}

\subsection{Cost ledger}

\Cref{tab:ledger} counts circuit executions on hardware (1,024 shots each) for training, model selection and
threshold calibration. We report three definitions, from the kit's to the most conservative:
\begin{enumerate}
\item \emph{Kit definition} (training and model selection): baseline 2,859,360 against MOCO 6,490, a reduction of
99.77\%.
\item \emph{With the noise meter} (\cref{sec:meter}): baseline 2,916,322 against MOCO $6{,}490+1{,}000=7{,}490$,
$-99.74\%$.
\item \emph{Most conservative}, threshold calibrated on the full noisy validation set (56,962 executions): MOCO
63,452, $-97.82\%$ (formula) and $-97.45\%$ against the Tracker count of the baseline.
\end{enumerate}
% ledger_final.csv (total executions; reduction vs baseline (formula) 0.978242, (Tracker) 0.974533);
% ledger_baseline.csv; ledger_moco_single_10.csv; noise_meter.csv (calibration executions 1000).
% Percentages 1 and 2 computed from those counts: 1 - 6490/2859360 = 99.77%, 1 - 7490/2916322 = 99.74%.

\begin{table}[t]
\centering\small
\caption{Cost ledger, circuit executions (formula; Tracker count in parentheses where it differs). Training +
model selection + threshold calibration.}
\label{tab:ledger}
% Source: results/task2/ledger_final.csv (and TABLES.md, section ledger_final.csv).
\setlength{\tabcolsep}{3pt}
\resizebox{\columnwidth}{!}{%
\begin{tabular}{lrrrr}
\toprule
Method & Train & Select & Calib. & Total\\
\midrule
Baseline (65/ex.) & 1,150,500 & 1,708,860 & 56,962 & 2,916,322\\
Baseline (Tracker, 41/ex.) & 725,700 & 1,708,860 & 56,962 & 2,491,522\\
R2 warm start (49/ex.) & 298,835 & 113,920 & 56,962 & 469,717\\
\textbf{MOCO} (2/row) & \textbf{6,490} & \textbf{0} & 56,962 & \textbf{63,452}\\
MOCO, noise meter & 6,490 & 0 & 1,000 & 7,490\\
\bottomrule
\end{tabular}}
\end{table}

\begin{figure}[t]
\centering
\includegraphics[width=\columnwidth]{\figdir/c_ledger.png}
\caption{Cost ledger (log scale). Source: \texttt{c\_ledger.png}, from \texttt{ledger\_final.csv}.}
\label{fig:ledger}
\end{figure}

\subsection{Finite shots and the shot tax}

The protocol evaluates exact expectation values. With $S=1024$ shots per transaction at $p=0.05$, test AUPRC is
0.705 for MOCO, 0.643 for R2 with warm start and 0.007 for the baseline, which falls to chance level
(\cref{fig:shottax}). The measured compression at $p=0.05$ is $0.871=\lambda^2$ for MOCO, 0.795 for R2 and 0.045
for the baseline, i.e.\ shot taxes of $1.32\times$, $1.58\times$ and at least $504\times$ (between $250\times$ and
$2.2\times10^5\times$ over seeds). For the baseline the tax is only a lower bound: its compression is not uniform
and its ranking is lost even with infinite shots.
% noise_shots.csv: shots = 1024, p = 0.05, mean over seeds: MOCO 0.705, R2 (moco) 0.643, baseline VQC 0.007.
% shot_tax.csv (p = 0.05): compression_mean MOCO 0.8711, R2 (moco) 0.7946, baseline 0.0445; shot_tax 1.3178,
% 1.5838, 503.94; baseline shot_tax_seed_min 250.06, seed_max 2.23e5.

\begin{figure}[t]
\centering
\includegraphics[width=\columnwidth]{\figdir/compact/e_shot_tax.png}
\caption{Shot tax $1/c^2$ against $p$ (diagnostic, measured on test, not used for any choice). Source:
\texttt{e\_shot\_tax.png}, from \texttt{shot\_tax.csv}.}
\label{fig:shottax}
\end{figure}

\subsection{All 492 frauds and the AUPRC/\texorpdfstring{$\rho$}{rho} paradox}

The official test set contains only 98 frauds. We therefore ran a stratified 5-fold cross-validation of the whole
dataset in which every fraud is tested once (\cref{fig:cv}). MOCO reaches $0.721\pm0.020$ at every $p$ and the
logistic regression on the same inputs $0.739\pm0.030$; the paired difference per fold is $-0.018\pm0.026$, so
the two cannot be distinguished, and the small advantage of MOCO on the official test set does not persist. The
baseline reaches $0.499\pm0.137$, with a spread over folds about seven times that of MOCO.
% cv/cv_per_fold.csv: AUPRC p = 0: MOCO 0.721 +- 0.020, logistic regression 1:10 0.739 +- 0.030, baseline
% 0.499 +- 0.137; paired MOCO - LR 1:10 per fold mean -0.018, std 0.026. 0.137 / 0.020 = 6.9.

The same folds show the paradox reported by \citet{dinut2026}: at low noise the baseline's AUPRC rises in 5 of 5
folds at $p=0.001$ and $p=0.01$ (mean $0.529\pm0.118$ at $p=0.01$), while its ranking decays ($\rho=0.931\pm0.122$
at $p=0.01$, $0.669\pm0.306$ at $p=0.05$). A single metric can thus report an improvement where the model gets
worse. Its drop at $p=0.05$ appears in only 3 of 5 folds. MOCO cannot show the paradox, because its ranking is
invariant.
% cv/cv_per_fold.csv: baseline VQC AUPRC rises over p = 0 in 5/5 folds at p = 0.001 and 0.01, 2/5 at 0.05;
% drops at 0.05 in 3/5; rho 0.931 +- 0.122 (0.01), 0.669 +- 0.306 (0.05).

\begin{figure}[t]
\centering
\includegraphics[width=\columnwidth]{\figdir/compact/i_cv_all_frauds.png}
\caption{Five-fold cross-validation on all 492 frauds. Source: \texttt{i\_cv\_all\_frauds.png}, from
\texttt{cv/cv\_per\_fold.csv}.}
\label{fig:cv}
\end{figure}

\subsection{Shallow entanglement with warm start (R2)}

R2 inserts three CNOTs between the encoding rotations and MOCO's trainable layer:
$RY(x)\to\mathrm{Rot}(\alpha)\to$ CNOT chain $\to\mathrm{Rot}(\beta)\to$ weighted $Z$ readout, 15 gates and 18
noise locations. With $\alpha=0$ it reduces to MOCO with a CNOT-chain feature map, whose last layer has a convex
global optimum; we use that optimum as a warm start for Adam. The only hypothesis stated before the run, that the
warm start beats a random initialization, holds: AUPRC $0.667\pm0.010$ against $0.635\pm0.018$, with half the
spread over seeds (\cref{tab:official}). Training inside the noise model at $p=0.05$ brings nothing
($0.666\pm0.009$). R2 stays below MOCO and costs 469,717 executions.
% circuits.csv: R2 entangled 15 gates, 18 noise locations. entangled_per_seed.csv: R2 (moco) 0.667 +- 0.010,
% R2 (random) 0.635 +- 0.018, R2 (moco-noisy) 0.666 +- 0.009 at p = 0. ledger_final.csv: 469717.

\subsection{Ablations}

\begin{table}[t]
\centering\small
\caption{Feature-map ablations with the same convex training (test AUPRC, 5 seeds).}
\label{tab:ablation}
% Sources: circuits.csv; moco_per_seed_{single,pair,linear_cnot,kit_zz}_10.csv (TABLES.md).
\setlength{\tabcolsep}{3pt}
\resizebox{\columnwidth}{!}{%
\begin{tabular}{lcccc}
\toprule
Map & CNOT & $p=0$ & $p=0.05$ & $\rho$ ($p=0.05$)\\
\midrule
single & 0 & $0.716\pm0.007$ & $0.716\pm0.007$ & 1.000\\
pair & 0 & $0.707\pm0.024$ & $0.707\pm0.024$ & 1.000\\
linear CNOT & 3 & $0.675\pm0.034$ & $0.667\pm0.032$ & $0.992\pm0.002$\\
kit ZZ map & 24 & $0.440\pm0.104$ & $0.450\pm0.114$ & $0.870\pm0.126$\\
\bottomrule
\end{tabular}}
\end{table}

\Cref{tab:ablation} changes only the feature map. Three CNOTs already lower AUPRC and break the invariance; a
second rotation per qubit (pair) adds nothing, which points to the information in four principal components as
the ceiling. With the kit's ZZ encoding, even the globally optimal readout stays at $0.440\pm0.104$: the encoding,
not the training, limits the baseline.

\subsection{Adaptive shot triage}

On hardware a transaction far from the threshold does not need 1,024 shots. One shot reads all four qubits and
gives $h=\sum_q w_q z_q\in[-1,1]$ with mean $\ex H_p$. We measure in batches of $b$ shots and stop as soon as an
anytime-valid confidence interval for $\ex H_p$ excludes the threshold $t_h=1-2t_p$, with a cap $S_{\max}$.
Three intervals were implemented: Hoeffding's \citep{hoeffding1963} and the empirical Bernstein bound of
\citet{maurer2009}, both with a union bound $\alpha_k=6\alpha/(\pi^2k^2)$ over batches, and the predictable
plug-in empirical-Bernstein confidence sequence of \citet{waudbysmith2024betting}. In all three cases the
probability that the interval misses the mean at any batch is at most $\alpha$, so the probability that a
transaction is stopped early with a decision different from that of its exact score is at most $\alpha$. In
coverage tests on 10,000 simulated sequences all three stayed below $\alpha$, whereas a fixed-sample interval
rechecked at every batch exceeded it (\texttt{tests/test\_adaptive.py}).

\paragraph{Selection on validation.} We chose the encoding scale $q$ (\cref{sec:robust}), the interval, $b$,
$\alpha$ and $S_{\max}$ on validation at $p=0.05$ (270 MOCO configurations, 5 seeds): highest validation F1,
ties within 0.005 broken by fewer shots. The winner is the kit encoding with the plug-in interval, $b=16$,
$\alpha=0.05$, $S_{\max}=4096$ (validation F1 0.788 at 208.4 mean shots). The interval is the main lever: at tied
F1, the Maurer--Pontil interval needs 991.7 shots and Hoeffding 917.9.
% adaptive_selection_on_validation.csv: MOCO kit, plugin-bernstein, b 16, alpha 0.05, cap 4096: F1 0.7877,
% mean shots 208.3775 (selected, headline); bernstein cheapest within 0.005 of its best F1: 991.7; hoeffding 917.9.

\paragraph{Test results.} At $p=0.05$ the triage uses $208.5\pm46.8$ shots per transaction instead of 1,024
($-80\%$) at the same F1 ($0.801\pm0.018$ against $0.797\pm0.019$ with 1,024 fixed shots), with zero wrong early
decisions in every MOCO run. At the same fixed budget for everyone (209 shots) F1 drops to $0.702\pm0.060$. Frauds,
which lie closer to the threshold, use more shots than legitimate transactions ($739.5$ against $207.6$). The
mean shots grow like $\lambda^{-4}$: $161.0\to208.5\to579.6$ from $p=0$ to $0.05$ and $0.2$, ratios $1.30$ and
$3.60$ against $\lambda^{-4}=1.32$ and $3.46$. AUPRC computed from the scores at stopping time drops from 0.716 to
0.700, because the triage deliberately leaves the far legitimate tail coarsely ranked. The baseline cannot be
triaged under noise: at $p=0.05$, 99.7\% of its transactions reach the cap and F1 stays at $0.091\pm0.072$.
% adaptive_shots.csv (cap 4096): MOCO kit adaptive p = 0.05: mean shots 208.489 +- 46.819, shots fraud 739.461,
% shots legit 207.574, F1 0.801 +- 0.018, AUPRC 0.700, wrong early decisions 0; fixed 1024 F1 0.797 +- 0.019;
% fixed same budget S = 209, F1 0.702 +- 0.060; p = 0: 160.982; p = 0.2: 579.576; exact AUPRC 0.716;
% baseline VQC adaptive p = 0.05: capped 0.997, F1 0.091 +- 0.072. shot_tax.csv: 1.3178, 3.4578.
% adaptive_checks.csv: 151 checks, all PASS.

\paragraph{Shots, not executions.} Each batch is one execution, so the triage costs 13.0 executions per
transaction at $p=0.05$ instead of one; the execution counts of \cref{tab:ledger} are unchanged. A
\texttt{qml.Tracker} run on the kit's noisy device (5 rows) counted 79 executions and 1,264 shots, equal to the
ledger formula $\sum_iK_i$ and $b\sum_iK_i$.
% ledger_adaptive.csv: MOCO kit adaptive p = 0.05: inference executions per transaction 13.031; Tracker check
% "executions 79 = formula 79, shots 1264 = formula 1264".

\paragraph{Confidence sweep.} On validation only, with the selected configuration fixed, lowering the confidence
$1-\alpha$ from 0.999 to 0.51 reduces the mean shots from 404.7 to 94.8 at $p=0.05$ ($4.3\times$); F1 stays within
0.01 of its value at $\alpha=0.05$ down to a confidence of 0.7 and loses about 0.02 at 0.51, and the rate of wrong
early decisions never exceeds 0.0002 (\cref{fig:sweep}). $\alpha=0.05$ was fixed before the sweep and is not changed.
% confidence_sweep.csv: p = 0.05 mean shots 404.7075 (0.999), 94.8034 (0.51); F1 0.7844 (0.95), 0.7865 (0.7),
% 0.7615 (0.51); max wrong early decisions 0.000176.

\begin{figure}[t]
\centering
\includegraphics[width=\columnwidth]{\figdir/l_confidence_sweep.png}
\caption{Confidence sweep of the adaptive triage on validation (5 seeds). Source:
\texttt{l\_confidence\_sweep.png}, from \texttt{confidence\_sweep.csv}.}
\label{fig:sweep}
\end{figure}

\subsection{Noise meter}\label{sec:meter}

Because $s_0$ is computable classically and $\mathbb E[\hat s_p]-\tfrac12=\lambda^2(s_0-\tfrac12)$ exactly, the
least-squares slope of $(\hat s_p-\tfrac12)$ on $(s_0-\tfrac12)$ over $N_{\mathrm{cal}}$ unlabeled calibration
transactions is an unbiased estimate of $\lambda^2$, and $\hat p=\tfrac34(1-\sqrt{\hat\lambda^2})$. With $S=1024$
shots per calibration transaction, the relative RMSE of $\hat p$ for $N_{\mathrm{cal}}=50/200/1000$ is
47\%/26\%/11\% at $p=0.01$, 11\%/5.3\%/2.4\% at $p=0.05$ and 3.7\%/1.8\%/0.8\% at $p=0.2$, falling roughly as
$1/\sqrt{N_{\mathrm{cal}}S}$. With $t_{\hat p}=g_{\hat p}(t_0)$, test F1 stays within 0.0003 of the F1 of the
threshold calibrated on the full noisy validation set at every $N_{\mathrm{cal}}$. Threshold calibration thus costs
1,000 executions instead of 56,962. The rule fixed before the run (relative error at most 10\% at every $p$) is
not met at $p=0.01$, where the error is 11\%, without effect on F1.
% noise_meter.csv: MOCO kit, least squares, S = 1024, mean over seeds of relative RMSE: p = 0.01: 0.4713,
% 0.2567, 0.1119; p = 0.05: 0.1068, 0.0528, 0.0244; p = 0.2: 0.0368, 0.0181, 0.0081; max |F1 drop| 0.0003.

\subsection{Robust encoding (unselected comparator)}\label{sec:robust}

The kit's angle scaling is fixed by the full range of each principal component, which outlying frauds set. A
robust variant fits the scale on the $[q,100-q]$ percentiles of the training subset and clips to $[0,\pi]$. On
validation it cuts the triage cost further, to 91.6 shots at $q=0.5\%$ and 33.8 at $q=5\%$, but lowers the
validation F1 (0.778 and 0.732 against 0.788), so the selection rule kept the kit encoding. Carried to test as a
comparator only, MOCO with $q=0.5\%$ reaches AUPRC $0.738\pm0.006$ at every $p$, and the logistic regression on the
same robust inputs $0.731\pm0.004$ (\cref{tab:official}); the higher test AUPRC does not change a choice made on
validation.
% adaptive_selection_on_validation.csv: plugin-bernstein, cheapest config within 0.005 of each model's best F1:
% q = 0.5%: 91.6 shots, F1 0.7775; q = 5%: 33.8 shots, F1 0.7323. moco_per_seed_single_10_q0.5.csv.

\section{Limitations}

\begin{itemize}
\item \textbf{Classically simulable.} The final state is a product state; scores are closed-form functions of the
inputs. As every four-qubit circuit in this challenge, MOCO is simulable, and we claim no quantum advantage
\citep{bowles2024benchmark,cerezo2023simulability}.
\item \textbf{Equivalent to a logistic regression} on $(\sin x_q,\cos x_q)$. MOCO matches, and does not beat, the
logistic regression on the same inputs, on the official test set and on all 492 frauds.
\item \textbf{Depolarizing noise only.} The theorem relies on the unital, isotropic channel of the official
model; amplitude damping or readout errors would add bias terms and need separate analysis.
\item \textbf{Four PCA inputs.} The precision-recall ceiling comes from the information in four principal
components: at 80\% recall the test precision is $0.817\pm0.009$, at 90\% only $0.072\pm0.032$.
% precision_at_recall.csv: max precision at recall 0.8: 0.817 +- 0.009; 0.9: 0.072 +- 0.032.
\item \textbf{98 test frauds.} Test metrics move by whole frauds; we therefore add the cross-validation on all 492.
\item \textbf{Triage cuts shots, not executions}, and its AUPRC is lower than with exact scores.
\end{itemize}

\section{Related work}

\emph{Data re-uploading and encoding.} Re-uploading circuits \citep{perezsalinas2020reuploading} and the
analysis of encodings as Fourier series \citep{schuld2021encoding} show that the encoding fixes the function class
of a quantum model; MOCO is the extreme case of a single upload per qubit, whose class is linear in
$(\sin x,\cos x)$. The ZZ feature map \citep{havlicek2019} used by the baseline is the opposite choice.
\emph{Classical surrogates and shadows.} Classical surrogates of quantum models \citep{schreiber2023surrogates} and
classical shadows \citep{huang2020shadows} learn from data measured once; MOCO does the same and compiles the
result back into the circuit. \emph{Simulability and trainability.} Barren plateaus \citep{mcclean2018barren},
noise-induced plateaus \citep{wang2021noise} and initialization strategies \citep{grant2019init} motivate our
convex training and the warm start of R2; the observation that provably trainable models are often classically
simulable \citep{cerezo2023simulability} applies to MOCO. \emph{Benchmarking.} \citet{bowles2024benchmark} argue
that quantum models must be compared with classical baselines on the same inputs; we follow this throughout.

\section{Conclusion}

MOCO turns the training of a variational classifier into a convex fit on data measured once and compiles the
optimum back into an 8-gate product circuit. Under the official depolarizing model its ranking, and with the
analytically transferred threshold every threshold metric, are invariant at every noise level: test AUPRC
$0.716\pm0.007$ against the baseline's drop to $0.405\pm0.177$, for 6,490 training executions instead of 2,859,360.
Noise remains a cost on hardware, a shot tax of $\lambda^{-4}$, which adaptive triage and a noise meter reduce. The
model matches a logistic regression on the same inputs; the gain is robustness, stability and cost. Next steps are
hardware noise beyond depolarizing channels and more qubits, which cost nothing in robustness for product
circuits.
% Numbers: moco_per_seed_single_10.csv, baseline_per_seed.csv, ledger_moco_single_10.csv, ledger_baseline.csv.

\paragraph{Reproducibility.} \texttt{make task2-test} runs the protocol, physics and interval tests;
\texttt{make task2-experiments} reruns every experiment with fixed seeds; \texttt{make task2-light} regenerates
all tables, figures and slides from the CSVs.

\bibliographystyle{unsrtnat}
\bibliography{refs_verified}% switch to refs once the team's paper/refs.bib is pushed

\end{document}
